\documentclass[12pt, a4paper]{article}

\usepackage[spanish]{babel}
\usepackage[utf8]{inputenc}
\usepackage[T1]{fontenc}
\usepackage{setspace}
\usepackage[margin=1in]{geometry}
\usepackage{amsmath}
\usepackage{amssymb}
\usepackage{graphicx}
\usepackage{hyperref}
\usepackage{booktabs}
\usepackage{array}

\onehalfspacing
\renewcommand{\arraystretch}{1.2}

\title{Propensión al Consumo de Marihuana en Colombia:\\ Benchmark MCO y Probit con Datos de Encuesta}

\author{Santiago Tupaz$^{1,2}$ \and Moisés Quintero$^{1}$ \and Vanessa Rodrigues$^{1}$ \\[0.5cm]
{\small $^{1}$ Universidad EAFIT, Programa de Economía, Medellín, Colombia} \\
{\small $^{2}$ Universidad Nacional de Colombia, Programa de Estadística, Sede Medellín} \\[1cm]
{\small \textbf{Profesora Supervisora:} Paula María Almonacid Hurtado, Universidad EAFIT}}

\date{\small 10 de marzo de 2026}

\begin{document}

\begin{titlepage}
    \centering
    \vspace*{3cm}

    {\Large \textbf{Universidad EAFIT}}\\[0.3cm]
    {\normalsize Programa de Economía}\\
    {\normalsize Ciencia de Datos}\\[2cm]

    {\large \textbf{Propensión al Consumo de Marihuana en Colombia}}\\[0.5cm]
    {\small Benchmark MCO y Probit con Datos de Encuesta}\\[2cm]

    {\normalsize
    \textbf{Autores:}\\
    Santiago Tupaz$^{1,2}$\\
    Moisés Quintero$^{1}$\\
    Vanessa Rodrigues$^{1}$
    }\\[1.5cm]

    {\small
    $^{1}$ Universidad EAFIT, Programa de Economía, Medellín, Colombia\\[0.3cm]
    $^{2}$ Universidad Nacional de Colombia, Programa de Estadística, Sede Medellín
    }\\[1.5cm]

    {\normalsize
    \textbf{Profesora Supervisora:}\\
    Paula María Almonacid Hurtado\\
    {\small Universidad EAFIT}
    }\\[2.5cm]

    {\normalsize
    Medellín, Colombia\\
    10 de marzo de 2026
    }

    \vfill
\end{titlepage}

\begin{abstract}
Este documento presenta la primera fase empírica del proyecto sobre consumo de marihuana en Colombia. El objetivo actual es modelar la propensión a reportar consumo en los últimos 12 meses utilizando información de la Encuesta Nacional de Consumo de Sustancias Psicoactivas y un conjunto básico de controles sociodemográficos. Para ello se construye una base de trabajo unificada a partir de una base limpia de consumo, el registro de personas y el capítulo sociodemográfico D2. Sobre esa base se estiman cuatro especificaciones de benchmark: dos modelos lineales de probabilidad y dos modelos Probit, con y sin información de precio. Los resultados muestran que la restricción por disponibilidad de precio reduce fuertemente la muestra útil, por lo que el coeficiente asociado al precio debe interpretarse solo como evidencia exploratoria. Esta fase deja establecido un benchmark econométrico sólido y una estructura reproducible para las siguientes etapas del proyecto, incluyendo la comparación con modelos de machine learning.

\textit{Palabras clave}: marihuana, propensión al consumo, MCO, Probit, benchmark econométrico, Colombia
\end{abstract}

\newpage

\tableofcontents

\newpage

\section{Introducción}

Este trabajo estudia la propensión individual a consumir marihuana en Colombia a partir de información de encuesta. La motivación central del proyecto es construir primero una base empírica clara y reproducible antes de pasar a modelos más complejos y a escenarios de política. En esta fase, el foco está en responder una pregunta simple y operativa: qué factores observables se asocian con la probabilidad de reportar consumo de marihuana en los últimos 12 meses.

La primera etapa se organiza alrededor de un benchmark econométrico. En vez de empezar con modelos complejos, se prioriza una estrategia escalonada: construcción de variables, limpieza, análisis exploratorio, estimación de un modelo MCO como referencia y estimación posterior de un modelo Probit para la misma variable dependiente. Esta secuencia permite evaluar calidad de datos, estabilidad de la muestra y sensibilidad de los coeficientes antes de avanzar a la fase de machine learning.

\section{Pregunta de investigación}

La pregunta de trabajo para esta fase es la siguiente:

\begin{quote}
\textbf{¿Qué factores observables se asocian con la probabilidad de consumir marihuana en los últimos 12 meses en Colombia?}
\end{quote}

De manera complementaria, se explora qué ocurre cuando se incorpora información de precio en la especificación. Sin embargo, dado que esa variable está disponible solo para una submuestra reducida, el análisis con precio debe entenderse como una extensión exploratoria y no como una estimación causal concluyente.

\section{Datos y construcción de la base}

\subsection{Fuentes de información}

La base final de modelado se construyó integrando tres fuentes principales:

\begin{table}[h!]
\centering
\small
\begin{tabular}{p{4.2cm} p{6.1cm} p{4.2cm}}
\toprule
\textbf{Archivo} & \textbf{Contenido} & \textbf{Llaves} \\
\midrule
\parbox[t]{4.2cm}{\texttt{base\_consumo\_drogas}\\\texttt{\_colombia\_limpia.xlsx}}
& Variables derivadas de consumo, gasto, precio, acceso y actitudes
& \texttt{id\_hogar}, \texttt{id\_encuesta}, \texttt{id\_persona}, \texttt{orden\_persona} \\
\texttt{personas\_processed.csv} & Sexo, edad y variables básicas de identificación individual & \texttt{directorio}, \texttt{secuencia\_encuesta}, \texttt{secuencia\_p}, \texttt{orden} \\
\texttt{d2\_capitulos.csv} & Escolaridad y variables sociodemográficas complementarias & \texttt{directorio}, \texttt{secuencia\_encuesta}, \texttt{secuencia\_p}, \texttt{orden} \\
\bottomrule
\end{tabular}
\caption{Fuentes principales para la base de modelado}
\label{tab:fuentes}
\end{table}

La verificación de llaves mostró coincidencia completa entre la base de consumo y las tablas \texttt{personas} y \texttt{d2} para las 3982 observaciones disponibles en la base limpia de consumo. Esa coincidencia permite construir una base integrada sin pérdida adicional en la etapa de merge.

La base integrada conserva por tanto las 3982 observaciones iniciales. Sin embargo, la muestra final para las especificaciones econométricas sin precio se reduce a 3980 porque dos observaciones presentan \texttt{D2\_05 = 9} (\textit{No sabe / No informa}) y quedan fuera al construir la agrupación educativa en tres categorías estables. Esta pérdida es mínima y no altera la interpretación general de los resultados, pero se documenta explícitamente por transparencia y reproducibilidad.

\subsection{Variables priorizadas}

El inventario operativo de variables se construyó a partir del diccionario de datos de la encuesta. Para esta fase se priorizaron las variables estrictamente necesarias para el benchmark:

\begin{table}[h!]
\centering
\small
\begin{tabular}{p{2.6cm} p{1.8cm} p{6.4cm} p{3.4cm}}
\toprule
\textbf{Variable limpia} & \textbf{Código} & \textbf{Definición} & \textbf{Uso} \\
\midrule
\texttt{consumo\_12m} & \texttt{K\_03} & ¿Ha consumido marihuana en los últimos 12 meses? & Variable dependiente \\
\texttt{precio\_compra} & \texttt{K\_09}, \texttt{K\_09\_VALOR} & Conocimiento y valor numérico del precio de un cigarrillo o porro de marihuana & Regresor de interés \\
\texttt{sexo} & \texttt{sexo} & Sexo de la persona & Control demográfico \\
\texttt{edad} & \texttt{edad} & Edad en años cumplidos & Control demográfico \\
\texttt{d2\_05} & \texttt{D2\_05} & Nivel educativo más alto alcanzado & Control socioeconómico \\
\texttt{frecuencia}\\{\_consumo} & \texttt{K\_04} & Frecuencia de consumo en los últimos 12 meses & Variable auxiliar \\
\texttt{facil\_marihuana} & \texttt{G\_06\_A} & Facilidad para conseguir marihuana & Proxy de acceso \\
\bottomrule
\end{tabular}
\caption{Variables principales integradas desde el inventario operativo}
\label{tab:variables}
\end{table}

Además, se dejaron identificadas para una segunda etapa variables de intensidad, como \texttt{cantidad\_consumo} (\texttt{K\_08}) y \texttt{gasto\_valor} (\texttt{K\_07}), que pueden servir para un modelo posterior de consumo intensivo o gasto.

\subsection{Disponibilidad y calidad de datos}

La principal restricción empírica de esta etapa es la disponibilidad desigual de variables. La base completa contiene 3982 observaciones, pero la cobertura efectiva cambia de forma importante según la variable:

\begin{itemize}
    \item \texttt{consumo\_12m}: 1719 registros no nulos en la base de consumo derivada.
    \item \texttt{precio\_compra}: 823 registros no nulos en la base derivada original.
    \item \texttt{cantidad\_consumo}: 1223 registros no nulos.
    \item \texttt{gasto\_valor}: 1014 registros no nulos.
    \item \texttt{frecuencia\_consumo}: cobertura completa en las 3982 observaciones.
\end{itemize}

La variable de precio requiere una aclaración adicional. En la base derivada original, los 823 registros no nulos de \texttt{precio\_compra} se distribuyen entre 603 personas que reportan consumo en los últimos 12 meses y 220 personas que no lo reportan. Luego de depurar códigos especiales y conservar únicamente precios válidos y positivos, la variable usable para el modelo se reduce a 756 observaciones. Finalmente, al exigir además covariables demográficas completas y una categoría educativa válida, la muestra econométrica con precio queda en 755 casos. Esta secuencia muestra que la información de precio no está distribuida aleatoriamente respecto al comportamiento de consumo y que la muestra con precio es una submuestra altamente seleccionada.

En consecuencia, el análisis se organiza en dos niveles de muestra:

\begin{enumerate}
    \item una muestra amplia sin precio, útil para establecer un benchmark general;
    \item una muestra restringida con precio, útil para explorar sensibilidad, pero no para sostener conclusiones causales fuertes.
\end{enumerate}

\subsection{Reglas de limpieza y recodificación}

La construcción de la base final siguió estas reglas:

\begin{enumerate}
    \item unión de las tres fuentes mediante llaves equivalentes;
    \item recodificación binaria de \texttt{K\_03} como variable de propensión a consumir;
    \item limpieza de variables numéricas positivas, eliminando códigos especiales como 8, 9, 98, 99, 998 y 999 cuando aparecen como no respuesta;
    \item transformación logarítmica del precio cuando la variable es válida y positiva;
    \item agrupación de educación en tres categorías estables: \textit{baja}, \textit{media} y \textit{superior}.
\end{enumerate}

La recodificación de la variable dependiente requiere un supuesto explícito. En el cuestionario, \texttt{K\_03} distingue al menos tres respuestas relevantes: 1 (\textit{Sí}), 2 (\textit{No}) y 9 (\textit{No contesta}). En la construcción de \texttt{propension\_12m}, el código 1 se transforma en 1, mientras que el código 2, el código 9 y los faltantes se tratan como 0. En consecuencia, la no respuesta se interpreta de la misma forma que la ausencia de consumo reportado. Este es un supuesto operativo útil para construir el benchmark, pero debe entenderse como una decisión metodológica provisional: si la no respuesta estuviera correlacionada con características observables o no observables de los individuos, podría introducir sesgo en etapas posteriores del análisis.

\section{Estrategia econométrica}

\subsection{Modelo benchmark MCO}

El primer benchmark se estimó con un modelo lineal de probabilidad:

\[
    y_i = \beta_0 + \beta_1 \log(precio_i) + \beta_2 edad_i + \beta_3 edad_i^2 + \beta_4 mujer_i + \beta_5 educacion_i + u_i
\]

donde \(y_i\) toma valor 1 si la persona reporta consumo de marihuana en los últimos 12 meses y 0 en caso contrario. Este modelo se usa como benchmark por su interpretación directa y por su utilidad para comparar especificaciones.

\subsection{Modelo Probit}

Como segunda especificación se estimó un modelo Probit con la misma estructura de covariables:

\[
    Pr(y_i = 1 \mid X_i) = \Phi(\alpha_0 + \alpha_1 \log(precio_i) + \alpha_2 edad_i + \alpha_3 edad_i^2 + \alpha_4 mujer_i + \alpha_5 educacion_i)
\]

En el reporte comparativo, las columnas del Probit se presentan mediante efectos marginales promedio para que la comparación con el MCO sea más clara.

\subsection{Especificaciones estimadas}

El análisis se estructura sobre dos muestras distintas derivadas de la misma base de modelado.

\begin{enumerate}
    \item \textbf{Muestra amplia} (\(n=3980\)): incluye todas las observaciones con edad, sexo y educación agrupada disponibles, sin requerir información de precio.
    \item \textbf{Muestra restringida} (\(n=755\)): corresponde al subconjunto de la muestra amplia para el cual existe además información válida de precio y covariables completas.
\end{enumerate}

Sobre cada muestra se estiman dos modelos, lo que genera cuatro especificaciones en total:

\begin{enumerate}
    \item \textbf{MCO completo}: muestra amplia sin precio;
    \item \textbf{MCO con precio}: muestra restringida con \(\log(precio)\);
    \item \textbf{Probit completo}: misma muestra amplia sin precio;
    \item \textbf{Probit con precio}: misma muestra restringida con \(\log(precio)\).
\end{enumerate}

Esta estructura hace explícita la lógica de dos muestras: por un lado, la asociación general entre variables sociodemográficas y propensión al consumo en la muestra más amplia posible; por otro, la asociación exploratoria entre precio y propensión en una submuestra seleccionada. La comparación entre ambas muestras es importante porque permite separar efectos de especificación de efectos derivados de la pérdida de observaciones.

\section{Resultados de benchmark}

\subsection{Tamaño muestral y ajuste}

La Tabla \ref{tab:fit} resume el número de observaciones y la métrica de ajuste para cada especificación.

\begin{table}[h!]
\centering
\caption{Resumen de muestra y ajuste}
\label{tab:fit}
\begin{tabular}{lrlrl}
\toprule
Modelo & N & Metrica & Valor & Incluye precio \\
\midrule
MCO completa & 3980 & R2 & 0.091 & No \\
MCO con precio & 755 & R2 & 0.061 & Si \\
Probit completo & 3980 & Pseudo R2 & 0.071 & No \\
Probit con precio & 755 & Pseudo R2 & 0.050 & Si \\
\bottomrule
\end{tabular}

\end{table}

Los modelos sin precio utilizan 3980 observaciones, mientras que los modelos con precio se reducen a 755. Esta caída de muestra es el principal límite empírico del análisis actual. En otras palabras, el benchmark sin precio describe mejor el patrón general de la muestra, mientras que el benchmark con precio debe leerse como una evidencia mucho más condicionada por selección muestral.

\begin{table}[h!]
\centering
\small
\caption{Cobertura de variables en la base de modelado}
\label{tab:cobertura}
\begin{tabular}{lrr}
\toprule
Variable & No nulos & \% cobertura sobre 3982 \\
\midrule
Consumo 12m & 1719 & 43.2 \\
Precio en base derivada & 823 & 20.7 \\
Precio válido para modelado & 756 & 19.0 \\
Edad & 3982 & 100.0 \\
Sexo & 3982 & 100.0 \\
Educación agrupable & 3980 & 99.9 \\
\bottomrule
\end{tabular}
\end{table}

\subsection{Coeficientes y efectos marginales}

La Tabla \ref{tab:benchmark} presenta los coeficientes del MCO y, para el Probit, los efectos marginales promedio.

\begin{table}[h!]
\centering
\small
\caption{Resultados comparados de benchmark}
\label{tab:benchmark}
\begin{tabular}{lllll}
\toprule
Variable & MCO completa & MCO con precio & Probit completo & Probit con precio \\
\midrule
Log precio &  & 0.044 (0.012) &  & 0.044 (0.012) \\
Mujer & -0.065 (0.016) & 0.006 (0.041) & -0.064 (0.016) & 0.004 (0.039) \\
Edad & -0.016 (0.003) & 0.013 (0.008) & -0.012 (0.004) & 0.010 (0.008) \\
Edad$^2$ & 0.000 (0.000) & -0.000 (0.000) & 0.000 (0.000) & -0.000 (0.000) \\
Educacion media & -0.007 (0.025) & 0.105 (0.055) & -0.005 (0.027) & 0.099 (0.049) \\
Educacion superior & -0.023 (0.025) & 0.079 (0.057) & -0.023 (0.027) & 0.074 (0.051) \\
\bottomrule
\end{tabular}

\begin{flushleft}\footnotesize Nota: la categoría de referencia para educación es Baja. Las columnas Probit reportan efectos marginales promedio. Errores estándar entre paréntesis. La variable dependiente es una medida binaria de propensión a consumir marihuana en los últimos 12 meses.\end{flushleft}

\end{table}

En esta comparación, la categoría de referencia para educación es \textit{baja}, que agrupa \textit{Ninguno}, \textit{Preescolar} y \textit{Básica primaria}. Por lo tanto, los coeficientes reportados para \textit{Educación media} y \textit{Educación superior} deben leerse en relación con ese grupo base.

\subsection{Lectura preliminar}

Los hallazgos más importantes de esta primera fase son los siguientes:

\begin{enumerate}
    \item En la muestra completa, la variable \textit{Mujer} presenta un coeficiente cercano a \(-0.065\) en MCO y un efecto marginal de \(-0.064\) en Probit. Bajo esta especificación, las mujeres muestran una menor propensión observada a reportar consumo de marihuana en los últimos 12 meses frente a los hombres.
    \item La edad aparece con signo negativo en la muestra completa tanto en MCO como en Probit, lo que sugiere una menor propensión observada al consumo a medida que aumenta la edad.
    \item En la muestra restringida con precio, el coeficiente de \(\log(precio)\) es positivo y estadísticamente distinto de cero en ambos modelos, con magnitud cercana a \(0.044\). Este resultado no debe interpretarse como elasticidad causal. La lectura correcta en esta etapa es que, dentro de una submuestra específica y altamente seleccionada, el precio aparece correlacionado con la propensión observada.
    \item Los coeficientes de educación cambian entre la muestra completa y la muestra con precio. Esto refuerza la idea de que la pérdida de observaciones por disponibilidad de precio altera la composición de la muestra y, con ello, la estabilidad de las estimaciones.
\end{enumerate}

\section{Reproducibilidad de la fase}

La fase quedó organizada para que el análisis pueda repetirse de forma transparente:

\begin{itemize}
    \item el inventario operativo de variables quedó documentado en \texttt{data/processed/variables.tex};
    \item la base de modelado se genera desde funciones en \texttt{src/cannabis\_tax/analysis/modeling.py};
    \item el análisis exploratorio y el benchmark se ejecutan desde el notebook \texttt{notebooks/01\_eda\_y\_mco\_propension.ipynb};
    \item las tablas del paper se exportan automáticamente a \texttt{reports/tables/}.
\end{itemize}

Esto deja la primera fase lista para revisión y facilita la transición a la siguiente etapa del proyecto.

\section{Conclusiones y siguiente fase}

La primera fase ya permite sostener varias conclusiones operativas. Primero, la pregunta de propensión al consumo en los últimos 12 meses está correctamente aterrizada en una variable observada y documentada del capítulo K de la encuesta. Segundo, la integración entre la base limpia de consumo y los módulos sociodemográficos funciona sin pérdida por llaves. Tercero, el benchmark MCO y Probit ya está completamente implementado y ofrece un punto de comparación claro para modelos posteriores.

Al mismo tiempo, el trabajo deja identificada la principal limitación de la etapa: la baja cobertura de la variable de precio. Por eso, la fase siguiente no debería centrarse en refinar indefinidamente el benchmark, sino en avanzar hacia el componente de machine learning, manteniendo el benchmark econométrico como referencia y evaluando si el desempeño predictivo mejora con modelos más flexibles.

\end{document}